\chapter{Homocystinuria}
\index{Homocystinuria}
\label{sec:Homocystinuria}

\section{Overview}
Homocystinuria (classic form) is an autosomal recessive disorder caused by deficiency of cystathionine $\beta$-synthase (CBS), the enzyme that converts homocysteine to cystathionine in the transsulfuration pathway. The incidence is approximately 1 in 200,000--300,000 live births.
\section{Causes}
This condition happens because of accumulation of homocysteine and methionine, producing a characteristic clinical phenotype resembling Marfan syndrome.     
\section{Risk Factors}

\begin{itemize}[noitemsep]
  \item Genetic predisposition and family history
  \item Advancing age
  \item Lifestyle factors (diet, exercise, smoking, alcohol)
  \item Environmental exposures
  \item Underlying medical conditions
  \item Medication use
\end{itemize}
\section{Signs and Symptoms}

\subsection{Common symptoms}
\begin{itemize}[noitemsep]
  \item You may experience tall, thin build, ectopia lentis (downward lens dislocation), myopia, intellectual disability (variable), skeletal abnormalities (scoliosis, pectus excavatum, arachnodactyly, genu valgum), osteoporosis, and a high risk of thromboembolism (the most life-threatening feature). 
  \end{itemize}

\subsection{Emergency warning signs}
\begin{itemize}[noitemsep]
  \item Severe or worsening symptoms of homocystinuria require immediate medical evaluation
\end{itemize}
\section{Progression and Stages}
The condition may progress through different stages of severity over time. Early stages often involve mild or subtle symptoms that may be overlooked. Moderate stages involve more noticeable symptoms affecting daily function. Advanced stages may involve significant impairment and complications.
\section{Complications}
\begin{itemize}[noitemsep]
  \item Disease progression leading to worsening symptoms
  \item Secondary infections or injuries
  \item Side effects of treatment
  \item Reduced quality of life
  \item Emotional and psychological impact
\end{itemize}
\section{Diagnosis}
Comprehensive medical history and physical examination Laboratory tests to assess organ function and rule out other conditions Imaging studies to visualize affected structures Specialized testing depending on the affected system Consultation with relevant specialists Monitoring of disease progression over time

\begin{itemize}[noitemsep]
  \item Comprehensive medical history and physical examination
  \item Laboratory tests to assess organ function
  \item Imaging studies to visualize affected structures
  \item Specialized testing depending on the affected system
  \item Consultation with relevant specialists
\end{itemize}
\section{Prevention}
\begin{itemize}[noitemsep]
  \item Doctors diagnose this using elevated total plasma homocysteine (greater than 100 \textmu mol/L), elevated methionine, and low cystine
  \item Newborn screening enables early detection.
  \item Maintain a healthy lifestyle with balanced diet and exercise
  \item Avoid known risk factors
  \item Attend regular medical check-ups for early detection
  \item Follow recommended screening guidelines
\end{itemize}
\section{Treatment Options}
Treatment focuses on pyridoxine for responsive patients

\begin{itemize}[noitemsep]
  \item Pyridoxine-nonresponsive patients require methionine restriction, betaine supplementation, and B12/folate supplementation
  \item Antiplatelet therapy reduces thromboembolic risk.
  \item Medications to manage symptoms and address underlying mechanisms
  \item Lifestyle modifications and self-care measures
  \item Physical therapy and rehabilitation
  \item Surgical intervention when indicated
  \item Regular monitoring and follow-up care
\end{itemize}
\section{Living with It}
\begin{itemize}[noitemsep]
  \item Follow your treatment plan as prescribed
  \item Maintain regular follow-up with your healthcare provider
  \item Make necessary lifestyle adjustments
  \item Seek support from family, friends, or support groups
  \item Stay informed about your condition
\end{itemize}
\section{Outlook}
The outlook is generally positive with early treatment. The outlook varies depending on the specific condition, severity at diagnosis, and response to treatment. Early intervention generally leads to better outcomes. Many conditions can be effectively managed with appropriate treatment. Regular follow-up care is important for monitoring and treatment adjustment.
